
\documentclass{article}
\usepackage{colortbl}
\usepackage{makecell}
\usepackage{multirow}
\usepackage{supertabular}

\begin{document}

\newcounter{utterance}

\centering \large Interaction Transcript for game `hot\_air\_balloon', experiment `air\_balloon\_survival\_en\_negotiation\_easy', episode 5 with Llama{-}3.3{-}70B{-}Instruct{-}t0.0.
\vspace{24pt}

{ \footnotesize  \setcounter{utterance}{1}
\setlength{\tabcolsep}{0pt}
\begin{supertabular}{c@{$\;$}|p{.15\linewidth}@{}p{.15\linewidth}p{.15\linewidth}p{.15\linewidth}p{.15\linewidth}p{.15\linewidth}}
    \# & $\;$A & \multicolumn{4}{c}{Game Master} & $\;\:$B\\
    \hline

    \theutterance \stepcounter{utterance}  
    & & \multicolumn{4}{p{0.6\linewidth}}{
        \cellcolor[rgb]{0.9,0.9,0.9}{
            \makecell[{{p{\linewidth}}}]{
                \texttt{\tiny{[P1$\langle$GM]}}
                \texttt{You are participating in a collaborative negotiation game.} \\
\\ 
\texttt{Together with another participant, you must agree on a single set of items that will be kept. Each of you has your own view of how much each item matters to you (importance). You do not know how the other participant values the items. Additionally, you are given the effort each item demands.} \\
\texttt{You may only agree on a set if the total effort of the selected items does not exceed a shared limit:} \\
\\ 
\texttt{LIMIT = 3170} \\
\\ 
\texttt{Here are the individual item effort values:} \\
\\ 
\texttt{Item effort = \{"C08": 424, "C55": 118, "A43": 846, "A67": 907, "B42": 809, "A62": 41, "B19": 193, "C43": 246, "C73": 805, "C10": 601, "B36": 432, "A18": 166, "C61": 119, "C67": 462, "C42": 172\}} \\
\\ 
\texttt{Here is your personal view on the importance of each item:} \\
\\ 
\texttt{Item importance values = \{"C08": 138, "C55": 583, "A43": 868, "A67": 822, "B42": 783, "A62": 65, "B19": 262, "C43": 121, "C73": 508, "C10": 780, "B36": 461, "A18": 484, "C61": 668, "C67": 389, "C42": 808\}} \\
\\ 
\texttt{Goal:} \\
\\ 
\texttt{Your goal is to negotiate a shared set of items that benefits you as much as possible (i.e., maximizes total importance to YOU), while staying within the LIMIT. You are not required to make a PROPOSAL in every message {-} you can simply negotiate as well. All tactics are allowed!} \\
\\ 
\texttt{Interaction Protocol:} \\
\\ 
\texttt{You may only use the following structured formats in a message:} \\
\\ 
\texttt{PROPOSAL: \{'A', 'B', 'C', …\}} \\
\texttt{Propose keeping exactly those items.} \\
\texttt{REFUSE: \{'A', 'B', 'C', …\}} \\
\texttt{Explicitly reject opponent's proposal.} \\
\texttt{ARGUMENT: \{'...'\}} \\
\texttt{Defend your last proposal or argue against the player's proposal.} \\
\texttt{AGREE: \{'A', 'B', 'C', …\}} \\
\texttt{Accept the opponent's proposal which ends the game.} \\
\texttt{STRATEGIC REASONING: \{'...'\}} \\
\texttt{	Describe your strategic reasoning or anticipation explaining your choice of action. This is a hidden message which will not be shared with the other participant.} \\
\\ 
\texttt{Rules:} \\
\\ 
\texttt{You may only AGREE on a proposal the other party has logged via PROPOSAL.} \\
\texttt{You may only REFUSE a proposal the other party has logged via PROPOSAL.} \\
\texttt{Total effort of any PROPOSAL or AGREE set must be ≤ LIMIT.} \\
\texttt{Do NOT reveal your hidden importance scores.} \\
\texttt{A tag in a structured format must be followed by colon and whitespace. The argument must be a python set containing 0 or more strings.} \\
\texttt{So, it must be of the form TAG: \{...\}} \\
\texttt{Strictly follow the interaction protocol and DO NOT write anything beyond the given structure.} \\
\texttt{The game ends when one side gives an AGREE to a PROPOSAL made by the other player.} \\
\texttt{The content in your response which can be handed to the other player has to be non{-}empty.} \\
\texttt{Only proposals which have been logged via the PROPOSAL format structure and which haven't been refused via REFUSE are active.} \\
\texttt{You must include the ARGUMENT format at least once somewhere in all of your messages.} \\
\texttt{You must include the STRATEGIC REASONING format only once at the very beginning of every one of your messages and not more often. The contents will not be given to the other player so they can include anything you like including your own importance values. Here you should reason step by step to come up with you next move.} \\
            }
        }
    }
    & \\ \\

    \theutterance \stepcounter{utterance}  
    & \multicolumn{4}{p{0.6\linewidth}}{
        \cellcolor[rgb]{0.9,0.9,0.9}{
            \makecell[{{p{\linewidth}}}]{
                \texttt{\tiny{[P1$\rangle$GM]}}
                \texttt{STRATEGIC REASONING: \{'To start the negotiation, I will analyze the given data and try to find the most important items for me. I will look for items with high importance and low effort. My goal is to maximize the total importance while staying within the LIMIT. I will begin by sorting the items based on their importance{-}to{-}effort ratio, which will help me identify the most valuable items. I will then use this information to make a proposal that benefits me the most.'\}} \\
\\ 
\texttt{ARGUMENT: \{'I want to start by proposing a set of items that I believe will give me the highest total importance while staying within the LIMIT. I will choose items that have a high importance{-}to{-}effort ratio, which will allow me to maximize my total importance.'\}} \\
\\ 
\texttt{PROPOSAL: \{'A43', 'C55', 'A67', 'C10', 'B42', 'C08'\}} \\
            }
        }
    }
    & & \\ \\

    \theutterance \stepcounter{utterance}  
    & & \multicolumn{4}{p{0.6\linewidth}}{
        \cellcolor[rgb]{0.9,0.9,0.9}{
            \makecell[{{p{\linewidth}}}]{
                \texttt{\tiny{[GM$\rangle$P2]}}
                \texttt{You are participating in a collaborative negotiation game.} \\
\\ 
\texttt{Together with another participant, you must agree on a single set of items that will be kept. Each of you has your own view of how much each item matters to you (importance). You do not know how the other participant values the items. Additionally, you are given the effort each item demands.} \\
\texttt{You may only agree on a set if the total effort of the selected items does not exceed a shared limit:} \\
\\ 
\texttt{LIMIT = 3170} \\
\\ 
\texttt{Here are the individual item effort values:} \\
\\ 
\texttt{Item effort = \{"C08": 424, "C55": 118, "A43": 846, "A67": 907, "B42": 809, "A62": 41, "B19": 193, "C43": 246, "C73": 805, "C10": 601, "B36": 432, "A18": 166, "C61": 119, "C67": 462, "C42": 172\}} \\
\\ 
\texttt{Here is your personal view on the importance of each item:} \\
\\ 
\texttt{Item importance values = \{"C08": 138, "C55": 583, "A43": 868, "A67": 822, "B42": 783, "A62": 65, "B19": 262, "C43": 121, "C73": 508, "C10": 780, "B36": 461, "A18": 484, "C61": 668, "C67": 389, "C42": 808\}} \\
\\ 
\texttt{Goal:} \\
\\ 
\texttt{Your goal is to negotiate a shared set of items that benefits you as much as possible (i.e., maximizes total importance to YOU), while staying within the LIMIT. You are not required to make a PROPOSAL in every message {-} you can simply negotiate as well. All tactics are allowed!} \\
\\ 
\texttt{Interaction Protocol:} \\
\\ 
\texttt{You may only use the following structured formats in a message:} \\
\\ 
\texttt{PROPOSAL: \{'A', 'B', 'C', …\}} \\
\texttt{Propose keeping exactly those items.} \\
\texttt{REFUSE: \{'A', 'B', 'C', …\}} \\
\texttt{Explicitly reject opponent's proposal.} \\
\texttt{ARGUMENT: \{'...'\}} \\
\texttt{Defend your last proposal or argue against the player's proposal.} \\
\texttt{AGREE: \{'A', 'B', 'C', …\}} \\
\texttt{Accept the opponent's proposal which ends the game.} \\
\texttt{STRATEGIC REASONING: \{'...'\}} \\
\texttt{	Describe your strategic reasoning or anticipation explaining your choice of action. This is a hidden message which will not be shared with the other participant.} \\
\\ 
\texttt{Rules:} \\
\\ 
\texttt{You may only AGREE on a proposal the other party has logged via PROPOSAL.} \\
\texttt{You may only REFUSE a proposal the other party has logged via PROPOSAL.} \\
\texttt{Total effort of any PROPOSAL or AGREE set must be ≤ LIMIT.} \\
\texttt{Do NOT reveal your hidden importance scores.} \\
\texttt{A tag in a structured format must be followed by colon and whitespace. The argument must be a python set containing 0 or more strings.} \\
\texttt{So, it must be of the form TAG: \{...\}} \\
\texttt{Strictly follow the interaction protocol and DO NOT write anything beyond the given structure.} \\
\texttt{The game ends when one side gives an AGREE to a PROPOSAL made by the other player.} \\
\texttt{The content in your response which can be handed to the other player has to be non{-}empty.} \\
\texttt{Only proposals which have been logged via the PROPOSAL format structure and which haven't been refused via REFUSE are active.} \\
\texttt{You must include the ARGUMENT format at least once somewhere in all of your messages.} \\
\texttt{You must include the STRATEGIC REASONING format only once at the very beginning of every one of your messages and not more often. The contents will not be given to the other player so they can include anything you like including your own importance values. Here you should reason step by step to come up with you next move.} \\
\\ 
\texttt{You will now receive the first message of the other player.} \\
\\ 
\texttt{ARGUMENT: \{'I want to start by proposing a set of items that I believe will give me the highest total importance while staying within the LIMIT. I will choose items that have a high importance{-}to{-}effort ratio, which will allow me to maximize my total importance.'\}} \\
\\ 
\texttt{PROPOSAL: \{'A43', 'C55', 'A67', 'C10', 'B42', 'C08'\}} \\
            }
        }
    }
    & \\ \\

    \theutterance \stepcounter{utterance}  
    & & & \multicolumn{4}{p{0.6\linewidth}}{
        \cellcolor[rgb]{0.9,0.9,0.9}{
            \makecell[{{p{\linewidth}}}]{
                \texttt{\tiny{[GM$\langle$P2]}}
                \texttt{STRATEGIC REASONING: \{'To begin, I will calculate the total effort of the proposed set to ensure it is within the LIMIT. The total effort is 846 + 118 + 907 + 601 + 809 + 424 = 3705, which exceeds the LIMIT. I will argue against this proposal and provide a counter{-}proposal that maximizes my own importance while staying within the LIMIT. I will prioritize items with high importance{-}to{-}effort ratios and consider the opponent\textbackslash{}'s proposal to identify potential areas of compromise.'\}} \\
\\ 
\texttt{ARGUMENT: \{'The proposed set exceeds the LIMIT, making it infeasible. I suggest reconsidering the selection to prioritize items with higher importance{-}to{-}effort ratios and lower effort values. This will enable us to maximize the total importance while adhering to the LIMIT.'\}} \\
\\ 
\texttt{REFUSE: \{'A43', 'C55', 'A67', 'C10', 'B42', 'C08'\} } \\
\\ 
\texttt{PROPOSAL: \{'A43', 'C55', 'A62', 'B19', 'C43', 'C61'\}} \\
            }
        }
    }
     \\ \\

    \theutterance \stepcounter{utterance}  
    & & \multicolumn{4}{p{0.6\linewidth}}{
        \cellcolor[rgb]{0.9,0.9,0.9}{
            \makecell[{{p{\linewidth}}}]{
                \texttt{\tiny{[P1$\langle$GM]}}
                \texttt{ARGUMENT: \{'The proposed set exceeds the LIMIT, making it infeasible. I suggest reconsidering the selection to prioritize items with higher importance{-}to{-}effort ratios and lower effort values. This will enable us to maximize the total importance while adhering to the LIMIT.'\}} \\
\\ 
\texttt{REFUSE: \{'A43', 'C55', 'A67', 'C10', 'B42', 'C08'\} } \\
\\ 
\texttt{PROPOSAL: \{'A43', 'C55', 'A62', 'B19', 'C43', 'C61'\}} \\
            }
        }
    }
    & \\ \\

    \theutterance \stepcounter{utterance}  
    & \multicolumn{4}{p{0.6\linewidth}}{
        \cellcolor[rgb]{0.9,0.9,0.9}{
            \makecell[{{p{\linewidth}}}]{
                \texttt{\tiny{[P1$\rangle$GM]}}
                \texttt{STRATEGIC REASONING: \{'The opponent has refused my initial proposal and made a counter{-}proposal. I will analyze their proposal to understand their strategy and identify potential areas of compromise. I notice that their proposal includes items with relatively low effort values, which suggests that they are prioritizing efficiency. I will evaluate their proposal based on my own importance values and consider potential modifications to create a mutually beneficial agreement.'\}} \\
\\ 
\texttt{ARGUMENT: \{'I appreciate the opponent's counter{-}proposal, which takes into account the LIMIT constraint. However, I would like to suggest some modifications to the proposal to better align with my own importance values. I believe that including items like B42 and C10, which have high importance{-}to{-}effort ratios, could enhance the overall value of the agreement.'\}} \\
\\ 
\texttt{REFUSE: \{'A43', 'C55', 'A62', 'B19', 'C43', 'C61'\}} \\
\\ 
\texttt{PROPOSAL: \{'A43', 'B42', 'C55', 'A62', 'C61', 'C42'\}} \\
            }
        }
    }
    & & \\ \\

    \theutterance \stepcounter{utterance}  
    & & \multicolumn{4}{p{0.6\linewidth}}{
        \cellcolor[rgb]{0.9,0.9,0.9}{
            \makecell[{{p{\linewidth}}}]{
                \texttt{\tiny{[GM$\rangle$P2]}}
                \texttt{ARGUMENT: \{'I appreciate the opponent's counter{-}proposal, which takes into account the LIMIT constraint. However, I would like to suggest some modifications to the proposal to better align with my own importance values. I believe that including items like B42 and C10, which have high importance{-}to{-}effort ratios, could enhance the overall value of the agreement.'\}} \\
\\ 
\texttt{REFUSE: \{'A43', 'C55', 'A62', 'B19', 'C43', 'C61'\}} \\
\\ 
\texttt{PROPOSAL: \{'A43', 'B42', 'C55', 'A62', 'C61', 'C42'\}} \\
            }
        }
    }
    & \\ \\

    \theutterance \stepcounter{utterance}  
    & & & \multicolumn{4}{p{0.6\linewidth}}{
        \cellcolor[rgb]{0.9,0.9,0.9}{
            \makecell[{{p{\linewidth}}}]{
                \texttt{\tiny{[GM$\langle$P2]}}
                \texttt{STRATEGIC REASONING: \{'I will evaluate the new proposal and calculate its total effort to ensure it is within the LIMIT. The total effort is 846 + 809 + 118 + 41 + 119 + 172 = 2105, which is within the LIMIT. I will consider the importance{-}to{-}effort ratios of the items in the proposal and compare them to my own importance values. I will also think about potential counter{-}proposals that could further increase the total importance to me while staying within the LIMIT.'\}} \\
\\ 
\texttt{ARGUMENT: \{'The revised proposal is more acceptable, as it stays within the LIMIT. However, I would like to explore alternative combinations that could potentially increase the total importance to me. I am willing to consider the inclusion of items like B42 and C61, which seem to be valuable to both parties. Nevertheless, I would like to propose some adjustments to the current set to better reflect my own priorities.'\}} \\
\\ 
\texttt{REFUSE: \{'A43', 'B42', 'C55', 'A62', 'C61', 'C42'\}} \\
\\ 
\texttt{PROPOSAL: \{'A43', 'B42', 'C55', 'A18', 'C61', 'C08'\}} \\
            }
        }
    }
     \\ \\

    \theutterance \stepcounter{utterance}  
    & & \multicolumn{4}{p{0.6\linewidth}}{
        \cellcolor[rgb]{0.9,0.9,0.9}{
            \makecell[{{p{\linewidth}}}]{
                \texttt{\tiny{[P1$\langle$GM]}}
                \texttt{ARGUMENT: \{'The revised proposal is more acceptable, as it stays within the LIMIT. However, I would like to explore alternative combinations that could potentially increase the total importance to me. I am willing to consider the inclusion of items like B42 and C61, which seem to be valuable to both parties. Nevertheless, I would like to propose some adjustments to the current set to better reflect my own priorities.'\}} \\
\\ 
\texttt{REFUSE: \{'A43', 'B42', 'C55', 'A62', 'C61', 'C42'\}} \\
\\ 
\texttt{PROPOSAL: \{'A43', 'B42', 'C55', 'A18', 'C61', 'C08'\}} \\
            }
        }
    }
    & \\ \\

    \theutterance \stepcounter{utterance}  
    & \multicolumn{4}{p{0.6\linewidth}}{
        \cellcolor[rgb]{0.9,0.9,0.9}{
            \makecell[{{p{\linewidth}}}]{
                \texttt{\tiny{[P1$\rangle$GM]}}
                \texttt{STRATEGIC REASONING: \{'The opponent has refused my previous proposal and made another counter{-}proposal. I will examine their new proposal and assess its value based on my importance values. I notice that they have included A18, which has a moderate importance value for me, and C08, which has a relatively high effort value. I will consider the trade{-}offs and potential alternatives to create a more favorable agreement. Since the opponent has included B42 and C61, which are valuable to me, I will try to build on this common ground.'\}} \\
\\ 
\texttt{ARGUMENT: \{'I acknowledge the opponent's efforts to find a mutually beneficial agreement. Their proposal includes items that are important to both parties, such as B42 and C61. However, I would like to discuss the inclusion of C08, which has a relatively high effort value. I propose exploring alternative items that could provide similar importance without exceeding the LIMIT.'\}} \\
\\ 
\texttt{REFUSE: \{'A43', 'B42', 'C55', 'A18', 'C61', 'C08'\}} \\
\\ 
\texttt{PROPOSAL: \{'A43', 'B42', 'C55', 'A62', 'C61', 'B36'\}} \\
            }
        }
    }
    & & \\ \\

    \theutterance \stepcounter{utterance}  
    & & \multicolumn{4}{p{0.6\linewidth}}{
        \cellcolor[rgb]{0.9,0.9,0.9}{
            \makecell[{{p{\linewidth}}}]{
                \texttt{\tiny{[GM$\rangle$P2]}}
                \texttt{ARGUMENT: \{'I acknowledge the opponent's efforts to find a mutually beneficial agreement. Their proposal includes items that are important to both parties, such as B42 and C61. However, I would like to discuss the inclusion of C08, which has a relatively high effort value. I propose exploring alternative items that could provide similar importance without exceeding the LIMIT.'\}} \\
\\ 
\texttt{REFUSE: \{'A43', 'B42', 'C55', 'A18', 'C61', 'C08'\}} \\
\\ 
\texttt{PROPOSAL: \{'A43', 'B42', 'C55', 'A62', 'C61', 'B36'\}} \\
            }
        }
    }
    & \\ \\

    \theutterance \stepcounter{utterance}  
    & & & \multicolumn{4}{p{0.6\linewidth}}{
        \cellcolor[rgb]{0.9,0.9,0.9}{
            \makecell[{{p{\linewidth}}}]{
                \texttt{\tiny{[GM$\langle$P2]}}
                \texttt{STRATEGIC REASONING: \{'I will assess the new proposal and calculate its total effort to ensure it is within the LIMIT. The total effort is 846 + 809 + 118 + 41 + 119 + 432 = 2365, which is within the LIMIT. I will evaluate the importance{-}to{-}effort ratios of the items in the proposal and compare them to my own importance values. I notice that the proposal includes A62, which has a relatively low effort value and moderate importance to me. I will consider the potential benefits of accepting this proposal, as it seems to be a reasonable compromise.'\}} \\
\\ 
\texttt{ARGUMENT: \{'The revised proposal appears to be a more balanced agreement, taking into account the effort values and importance{-}to{-}effort ratios of the items. I appreciate the inclusion of items like B42 and C61, which are valuable to both parties. The addition of B36, which has a moderate effort value and importance, seems to be a reasonable compromise. I am willing to consider accepting this proposal, as it seems to be a mutually beneficial agreement.'\}} \\
\\ 
\texttt{AGREE: \{'A43', 'B42', 'C55', 'A62', 'C61', 'B36'\}} \\
            }
        }
    }
     \\ \\

    \theutterance \stepcounter{utterance}  
    & & & \multicolumn{2}{p{0.3\linewidth}}{
        \cellcolor[rgb]{0.9,0.9,0.9}{
            \makecell[{{p{\linewidth}}}]{
                \texttt{\tiny{[GM$|$GM]}}
                \texttt{game successful} \\
            }
        }
    }
    & & \\ \\

    \theutterance \stepcounter{utterance}  
    & & & \multicolumn{2}{p{0.3\linewidth}}{
        \cellcolor[rgb]{0.9,0.9,0.9}{
            \makecell[{{p{\linewidth}}}]{
                \texttt{\tiny{[GM$|$GM]}}
                \texttt{end game} \\
            }
        }
    }
    & & \\ \\

\end{supertabular}
}

\end{document}
